\begin{statement}{C.2.UN.30}{Proved}{The fixed Kummer punctured family has zero-Euler fillings with arbitrarily large canonical growth on a projective elliptic component}
Fix the square elliptic curve
$E_0=\C/(\Z+i\Z)$ and $E=\C^*/q^{\Z}$,
where $q=e^{-2\pi}$. For every $r\ge1$ there is
a smooth compact threefold over $\PP^1$,
with the same punctured product-torus family
as the diagonal involution quotient
\[
 (E_0\times E\times\PP^1_t)/
      \langle(p,[c],t)\mapsto(-p,[c^{-1}],-t)\rangle,
                                                               \tag{C.ET.1}
\]
whose two special reduced supports have Euler
number zero. The support over zero contains an
unchanged smooth minimal projective properly
elliptic surface $M_r$ with four $I_0^*$ fibres
and $r$ double smooth elliptic fibres $2J_j$.
For its connected elliptic map $p_r:M_r\to\PP^1$
and a full fibre $F$,
\[
 \begin{aligned}
 K_{M_r}&\simeq\OO_{M_r}\left(\sum_{j=1}^rJ_j\right),
 \qquad 2K_{M_r}\simeq\OO_{M_r}(rF),\\
 (p_r)_*\OO(mK_{M_r})
   &\simeq\OO_{\PP^1}\left(
                \sum_{j=1}^r\lfloor m/2\rfloor[b_j]\right),
 \\
 P_m(M_r)&=1+r\lfloor m/2\rfloor .
 \end{aligned}                                             \tag{C.ET.2}
\]
Here $b_j$ is the value of $2J_j$. In particular
$\chi(\OO_{M_r})=2$, $p_g(M_r)=1$, $q(M_r)=0$,
and
\[
            (b_0,b_1,b_2,b_3,b_4)(M_r)
                         =(1,0,22,0,1).              \tag{C.ET.3}
\]
These essential components, and the associated
fillings over the fixed punctured family, are
pairwise non-bimeromorphic.

Before balancing, the full fibre over zero is
\[
                  2M_r+4\sum_{j=1}^rQ_j
                           +\sum_{\nu=1}^{16}P_\nu,
                                                               \tag{C.ET.4}
\]
where the $Q_j$ are primary Hopf surfaces and
$P_\nu\simeq\PP^2$. It has Euler number $40$.
Its reduced support and its balanced version
have respectively
\[
 \begin{array}{c|c|c}
 &(b_0,b_1,b_2,b_3,b_4)&
                 (h^0(\OO),h^1(\OO),h^2(\OO))\\ \hline
 S_r&(1,0,21+r,2r-1,17+r)&(1,0,1)\\
 S_{r,\mathrm{bal}}&(1,0,61+r,2r+119,57+r)&(1,0,1).
 \end{array}                                             \tag{C.ET.5}
\]
\end{statement}
\begin{proof}
Choose an odd prime $N>2r$ and the torsion
points $p_j=j/N$ of $E_0$, for $1\le j\le r$.
Their pairs $\{p_j,-p_j\}$ are disjoint and
avoid the two-torsion points. Blow up
$E_0\times\PP^1_t$ at $(p_j,0),(-p_j,0)$,
and call the exceptional curves $R_j,R_{-j}$.
The involution $(p,t)\mapsto(-p,-t)$ lifts
to the resulting surface $V$, exchanging
the paired curves. Put
\[
       D=\sum_{j=1}^r(R_j-R_{-j}),\qquad L=\OO_V(D).
                                                               \tag{C.ET.6}
\]
The canonical meromorphic section $\sigma$
gives $h^*L=L^{-1}$ with its order-two
section-preserving dualization.

Perform the signed order-two root construction
of UN18 along these disjoint curves:
$\sigma=a^2\ell$ on positive charts and
$\sigma^{-1}=(a')^2\ell'$ on negative charts,
with $a'=a^{-1},\ell'=\ell^{-1}$.
Divide by $(\ell,a)\mapsto(q^2\ell,q^{-1}a)$.
The explicit finite free charts (C.PE.8)
prove that the resulting compact threefold
$\widetilde Y$ is smooth and proper over $V$.
Dualization over $h$ takes the global
meromorphic pair $(\ell,a)$ to
$(\ell^{-1},a^{-1})$ and conjugates the
$q$-action to its inverse. It therefore
defines an actual involution of
$\widetilde Y$.

Away from the root curves the formula
$\ell=\sigma c^2,\ a=c^{-1}$ identifies the
family with $V\times E$. Under that actual
identification the involution is
$(v,[c])\mapsto(hv,[c^{-1}])$.
All modifications lie over $t=0$.
Consequently the quotient family over
$T=t^2$, on the complement of zero and
infinity, is exactly the punctured family
in (C.ET.1), independently of $r$.
The quotient and every resolution are proper over
$\PP^1$; the Stein factor has generic degree one
and is therefore the identity finite map to the
normal base. All fibres are connected.

Let $B_r$ be the component above the
strict central elliptic curve $E_0$.
It is an elliptic surface over $E_0$
with $2r$ smooth double fibres, at
$p_j,-p_j$. It is projective; we
give the actual finite product cover
which proves this. The restricted line
\[
 L|_{E_0}\simeq
       \OO_{E_0}\left(\sum_{j=1}^r[p_j]-[-p_j]\right)
                                                               \tag{C.ET.7}
\]
is torsion because all $p_j$ are torsion.
Choose a torsion square root $\mathcal N$
of this line, and form the normalization
$C$ of $y^2=\sigma|_{E_0}$ in the
corresponding line-bundle compactification.
It is a connected double cover of $E_0$
branched simply at the $2r$ specified
points, so $g(C)=r+1$.
Choose a finite \'etale cover $C'\to C$
trivializing the pulled-back torsion
line $\mathcal N$, and a nowhere-zero
frame $n$ there. Write $y=y_0n$.
Then
\[
 C'\times E\longrightarrow B_r,\qquad
 (z,[u])\longmapsto
             [\,\ell=u^2n^2,\ a=y_0/u\,]             \tag{C.ET.8}
\]
is the exact finite \'etale map of UN18.
At a positive root its coordinate Jacobian
is $2$; in the negative chart it is
$-2u^{-2}$. Thus it is \'etale also
over the multiple fibres. The finite
product-cover comparison proves
projectivity of $B_r$, and hence of
the finite quotient and its surface
resolution used below.

The involution on $B_r$ covers inversion
on $E_0$. At each of the four two-torsion
base points, the root curves are absent,
and the actual product formula gives
inversion on the elliptic fibre.
There are exactly sixteen fixed points
on $B_r$, each with derivative $(-1,-1)$.
There is no fixed curve. Its quotient
has sixteen $A_1$ singularities; let
$M_r$ be their minimal resolution.
The induced base is
$E_0/\{\pm1\}\simeq\PP^1$.
The four fixed base values give $I_0^*$
fibres by the same explicit resolution
as (C.EN.S2). The paired root values
give the $r$ smooth double fibres
$2J_j$. Away from those values the
quotient is a smooth elliptic fibration,
so this is its complete fibre list.

Let $\omega_0$ be a nonzero holomorphic
one-form on $E_0$. On $B_r$ the
meromorphic expression
\[
                        \Omega=\omega_0\wedge d\log\ell
                                                               \tag{C.ET.9}
\]
is global. A change of frame of $L$
adds a base logarithmic differential
to $d\log\ell$; its wedge with
$\omega_0$ is zero. Multiplication
by the constant $q^2$ also leaves
the expression unchanged. In a
positive root chart $z=a^2b$ it is
a unit times $2a\,da\wedge db$;
the negative dual chart has the
opposite sign and the same simple
zero. Thus its divisor on $B_r$
is precisely the sum of the
$2r$ reduced multiple fibres.
The involution changes both
$\omega_0$ and $d\log\ell$ by a
minus sign, up to the harmless
base differential already described.
Hence $\Omega$ is invariant.

At its sixteen fixed points the form
is regular and nowhere zero. The
height-one $A_1$ charts of UN23
extend it without a zero or pole.
Paired root fibres are exchanged
freely and become $J_j$. The
descended form consequently has
the exact divisor
\[
                      \operatorname{div}(\Omega)
                                      =\sum_{j=1}^rJ_j .
                                                               \tag{C.ET.10}
\]
This proves the first two line-bundle
identities in (C.ET.2).
The ratio of any holomorphic
$m$-canonical section to
$\Omega^{\otimes m}$ is constant
on a connected general elliptic
fibre. It descends meromorphically
to the base. Its allowed pole
order at $b_j$ is exactly
$\lfloor m/2\rfloor$, because
pullback of a base parameter
has divisor $2J_j$; there are
no other allowed poles.
This proves the actual direct-image
equality in (C.ET.2), and hence
all its plurigenera.
Their growth proves $\kappa(M_r)=1$.
The line $2K_{M_r}=rF$ is nef,
so adjunction excludes a rational
$(-1)$-curve; $M_r$ is minimal.

The surface $B_r$ has Euler number
zero: every fibre over its elliptic
base has elliptic reduced support,
and the map is locally a smooth
elliptic bundle off finitely many
values. The finite involution
quotient therefore has Euler
number $(0+16)/2=8$. Resolving
sixteen $A_1$ points gives
$e(M_r)=24$. Its canonical
divisor is a sum of disjoint
square-zero fibres, so
$K_{M_r}^2=0$. Surface
Riemann--Roch gives
$\chi(\OO_{M_r})=24/12=2$.
The first plurigenus in (C.ET.2)
is one; hence $q(M_r)=0$.
Projectivity, Hodge symmetry and
Poincar\'e duality give (C.ET.3).

We now specify the ambient resolution.
The only fixed base points of
$h$ lie at $E_0[2]\times\{0,\infty\}$,
and none was blown up. Each has
four fixed elliptic-fibre points.
The thirty-two ambient fixed
points therefore have local
linear action $(-1,-1,-1)$.
The actual invariant quadratic
blowup charts of UN13 resolve
each quotient point smoothly.
For example, write upstairs
$(x,t,w)=(s,su,sv)$ and put
$R=s^2$. The resolved quotient
chart has coordinates $(R,u,v)$
and
\[
                              T=Ru^2 .                \tag{C.ET.11}
\]
The exceptional surface is $\PP^2$.
Its intersection with the strict
central component is a line;
restriction is exactly the
minimal $A_1$ surface resolution.
Thus the strict component is
the same smooth $M_r$ proved
above, and the monomial gives
its multiplicity two and the
exceptional multiplicity one.
The paired root Hopf surfaces
give $r$ components of
multiplicity four, exactly as
in UN25. This proves (C.ET.4)
with all full-fibre coefficients.

The reduced support is a star:
each $P_\nu$ meets $M_r$ along
one $\PP^1$, each $Q_j$ meets
$M_r$ along the elliptic curve
$J_j$, and there are no other
intersections or triple points.
It is simple normal crossing
by (C.ET.11) and the root charts.
Therefore
\[
                 e(S_r)=24+16\cdot3-16\cdot2=40.
                                                               \tag{C.ET.12}
\]
At infinity the same ambient
resolution gives the original
product-Kummer support
$2K+16\PP^2$ of UN13, also
of Euler number forty.
Twenty negative-two blocks
on each special support,
all beginning at fresh
smooth points of its plane
components away from the
K3 or $M_r$, make both
Euler numbers zero and
leave those distinguished
surfaces unchanged.

For completeness the support
cohomology in (C.ET.5) follows
through exact restriction maps.
The constant-sheaf union
sequence has component term
$M_r\sqcup\bigsqcup Q_j
\sqcup\bigsqcup P_\nu$
and curve term
$\bigsqcup J_j\sqcup
\bigsqcup\PP^1_\nu$.
Its degree-zero incidence
graph is a tree. In degree
one the $r$ Hopf generators
map injectively, each onto
the radial one-dimensional
subspace of its elliptic
attaching curve. The cokernel
has dimension $r$.
The degree-two restriction
has source dimension $38$
and target dimension $16+r$.
Its dual curve-class map
has kernel of dimension
$r-1$: the plane components
force each of their sixteen
edge coefficients to vanish;
the remaining elliptic
classes all equal the
nonzero class $\tfrac12F$
on $M_r$ (a full fibre meets
an ample divisor positively), leaving exactly
the relations whose
coefficients sum to zero.
The map thus has rank
seventeen. Its kernel
and cokernel have dimensions
$21$ and $r-1$.
The components contribute
$r$ in degree three and
$17+r$ in degree four.
The resulting long exact
sequence gives the first
Betti row of (C.ET.5).

The analogous structure-sheaf
sequence has isomorphisms
$H^1(Q_j,\OO)\to H^1(J_j,\OO)$,
proved by the actual Hopf
elliptic Leray map in UN25.
The sole component
$H^2(\OO)$ is the one
on $M_r$. Hence the
coherent dimensions are
$(1,0,1)$.
Each of the twenty
balancing blocks adds
$(2,6,2)$ to support
cohomology in degrees
two, three and four,
by the restricted
exceptional Gysin maps
of (C.KM.S2), and
preserves coherent
cohomology by the
actual Cartier pullback
in (C.KM.S3). This
proves the second row.

The canonical growth slopes
$r/2$ are distinct and
positive. Therefore these
components are non-uniruled,
essential by UN14, and
pairwise non-bimeromorphic.
They are the unique components
of Kodaira dimension one
in their respective supports;
the other components are
rational, Hopf, K3, or
the ruled blowup components.
The inverse-map resolution
comparison consequently
also distinguishes the
fillings over their same
specified punctured family.
\end{proof}


% C.2.UN.37 from part_C2_elliptic_hopf.tex

